\documentclass[11pt,a4paper]{article}

\usepackage[utf8]{inputenc}
\usepackage[T1]{fontenc}
\usepackage{geometry}
\geometry{margin=1in}
\usepackage{microtype}
\usepackage{hyperref}
\hypersetup{colorlinks=true, linkcolor=blue, citecolor=blue, urlcolor=blue}
\usepackage{booktabs}
\usepackage{listings}
\usepackage{xcolor}
\usepackage{authblk}
\usepackage{amsmath}
\usepackage{enumitem}
\usepackage{natbib}
\usepackage{array}
\setlength{\emergencystretch}{3em}

\lstset{
  basicstyle=\ttfamily\footnotesize,
  breaklines=true,
  frame=single,
  framesep=4pt,
  rulecolor=\color{gray!30},
  showstringspaces=false,
  columns=fullflexible,
}

\title{RLM-FORGE\\
       Runtime-Lifted Recursive Language Models for Agent Infrastructure}

\author[1]{JQ Lee}
\affil[1]{Independent\\
\texttt{jqyu.lee@gmail.com}}

\date{May 3, 2026}

\begin{document}
\maketitle

\begin{abstract}
The immediate value of Recursive Language Models (RLMs) extends beyond
training a new recursive checkpoint. We argue that the RLM pattern can be
\emph{runtime-lifted} into an execution contract inside an existing agent
runtime, allowing ordinary provider-swappable LLM calls to inherit the
operational controls associated with recursion. RLM-FORGE implements this
claim by mapping RLM sub-calls onto the existing RPC mechanism of Hermes
Agent. Ouroboros owns recursive state, scheduling, trace replay, and evidence
validation. The result is a deployable runtime primitive for recursive
inference, not a new model architecture. The primitive preserves the
operational parts of RLM that matter most, including bounded context
partitioning, explicit sub-call graphs, replayable intermediate evidence,
and parent synthesis that can be rejected before state mutation. We introduce
TraceGuard, a deterministic evidence gate for parent synthesis, and evaluate
it with a claim-aware omitted-fact safety scorer. The live Hermes truncation
fixture shows score parity between vanilla and recursive execution after
scorer correction. A 24-cell live primary portability run across Hermes+GLM,
Claude Code, and Codex completes with 24 mandatory contract passes. The
systems claim is not that RLM-FORGE wins a benchmark. It is that RLM-FORGE
turns RLM from a research-time inference idea into a runtime surface that
existing agent infrastructure can execute, inspect, and enforce.
\end{abstract}

\section{Introduction}
\label{sec:intro}

Recursive Language Models, introduced by \citet{zhang2025rlm}, are usually
read as a new inference paradigm for long-context reasoning. That reading is
correct. It also hides a second and more immediately deployable consequence.
RLM is a runtime interface. If an LLM can be invoked on a bounded slice
of context, return structured evidence, and be called again under outer
control, then the runtime can supply the recursive substrate even when the
underlying model was not post-trained for recursion.

This paper takes that runtime-first interpretation seriously. We ask a
deliberately practical question. Can the advantages of RLM be moved out of a
research scaffold and into an existing agent runtime without requiring a new
model architecture, sandbox, or transport layer? RLM-FORGE answers yes for one
concrete stack. Hermes Agent supplies the callable inner execution layer.
Ouroboros supplies the recursion controller, trace store, and
acceptance-criteria tree. TraceGuard supplies evidence-gated parent synthesis.

The key distinction is between \emph{quality advantage} and \emph{control
advantage}. A single vanilla LLM call can be correct on a controlled fixture.
That does not erase the value of runtime recursion. The practical advantage
of RLM is that it changes the shape of inference. Context becomes
partitionable, work becomes compositional, intermediate claims become
inspectable, and parent synthesis can be blocked when it cites evidence the
children never produced. Those are runtime properties. They are useful even
before a broad live benchmark demonstrates a model-quality delta.

We make four concrete claims.

\begin{enumerate}[leftmargin=*]
\item \textbf{Runtime-lifted RLM.} The existing RPC pattern in Hermes maps
directly onto the RLM sub-call envelope. Hermes describes this pattern as
Python scripts calling tools through RPC to collapse multi-step pipelines into
zero-context turns~\citep{nous2025hermes}. This makes RLM executable as an
agent-runtime capability, not as a special model endpoint.

\item \textbf{Preserved RLM control surfaces.} RLM-FORGE preserves the
operational advantages of RLM at runtime. These include bounded context
windows, recursive sub-call graphs, explicit child evidence, and parent
synthesis under outer control.

\item \textbf{Incremental integration.} Two integration paths share the same
Hermes adapter without modifying the default control flow of the host system.
One path is a dedicated recursive scaffold (Section~\ref{sec:scaffold}).
The other is an existing acceptance-criterion decomposition pipeline
(Section~\ref{sec:decomposition}).

\item \textbf{Executable evidence discipline.} TraceGuard rejects structured
claims that lack accepted child evidence handles. In deterministic contract
ablations, evidence-gated Hermes-RLM has a 0.0 unsupported-claim rate while
an ungated RLM-shaped policy fails at 1.0. In the 24-cell live primary matrix
(Section~\ref{sec:experiment}), Hermes+GLM, Claude Code, and Codex all pass
the same 8-fixture RLM-FORGE+TraceGuard contract. The implementation also
includes a narrow repair loop for a missing-handle rejection class observed in
an earlier Hermes+GLM run. On the
original live truncation fixture, corrected scoring reports a tie rather than
the earlier false $+0.20$ recursive-RLM win.
\end{enumerate}

We do not claim novelty over the RLM paradigm itself, and we do not compare
to the natively recursive RLM-Qwen3-8B model post-trained by
\citet{zhang2025rlm}. Our claim is a systems claim. RLM can become a runtime
primitive that existing agent infrastructure can host, trace, and enforce.

\section{Background}
\label{sec:bg}

\paragraph{Recursive Language Models.}
\citet{zhang2025rlm} formalise an inference strategy in which an outer
LLM call invokes itself, or another LLM, on a transformed slice of
its input via a tool call $\mathrm{RLM}_M(\hat{q}, \hat{C})$. The original
paper realises the environment as a Python REPL that holds the input
context as a variable and exposes \texttt{peek}, \texttt{partition},
\texttt{transform}, and recursive call operations to the LM. The headline
result is that this strategy processes inputs up to two orders of
magnitude beyond model context windows while frequently outperforming
vanilla calls and common long-context scaffolds across four diverse
tasks.

\paragraph{Hermes Agent.}
Hermes Agent is an open-source self-improving AI agent runtime by Nous
Research. Two of its features matter for our purposes. First, Hermes is
provider-agnostic. A single CLI subcommand, \texttt{hermes model}, swaps
the inner LLM across hundreds of providers. Second, Hermes ships
an RPC tool mechanism described in its README as ``write Python scripts
that call tools via RPC, collapsing multi-step pipelines into
zero-context turns,'' along with isolated subagent spawn for parallel
work. The combination yields a runtime that is callable like a function
with structured JSON in and out, which is exactly what an RLM scaffold
needs as its recursive boundary.

\paragraph{Ouroboros.}
Ouroboros~\citep{ouroboros2026} is a workflow scaffold for AI coding
agents. It contributes two pieces relevant to our integration. They are a
typed acceptance-criteria tree (\texttt{ACTree}) with a depth cap of five
and a runtime adapter (\texttt{HermesCliRuntime}) that wraps the Hermes CLI
as a callable \texttt{AgentRuntime}. We use both unchanged.

\section{System Design}
\label{sec:system}

\subsection{Runtime lifting}
\label{sec:runtime-lifting}

We use \emph{runtime lifting} to mean the following transformation. An
inference pattern that is often described at the model or research
environment level is exposed as a stable runtime contract that ordinary agent
infrastructure can execute. For RLM, the contract is simple. The outer
scaffold selects a bounded context, invokes an inner LM on that context,
records the returned evidence, and decides whether to recurse, summarize,
synthesize, retry, or terminate.

This is stronger than wrapping a model in map-reduce. Map-reduce produces
chunk summaries and then trusts a reducer. Runtime-lifted RLM produces a
call graph. Each child call has a parent, a bounded input, a typed mode,
an evidence handle set, and an accepted or rejected result. The parent
synthesis step is then another runtime-validated operation rather than a
free-form final answer.

\begin{figure}[t]
\centering
\small
\setlength{\fboxsep}{6pt}
\begin{tabular}{c}
\fbox{\begin{minipage}{0.82\linewidth}
\centering
\textbf{Ouroboros controller}\\
select bounded context, schedule RLM nodes, own trace state
\end{minipage}}\\[0.5em]
$\downarrow$ \texttt{HermesCliRuntime} envelope\\[0.4em]
\fbox{\begin{minipage}{0.82\linewidth}
\centering
\textbf{Hermes bounded sub-calls}\\
execute local decomposition, chunk summary, or atomic work
\end{minipage}}\\[0.5em]
$\downarrow$ accepted child evidence manifest\\[0.4em]
\fbox{\begin{minipage}{0.82\linewidth}
\centering
\textbf{TraceGuard parent-synthesis gate}\\
validate every structured claim against accepted child evidence
\end{minipage}}\\[0.5em]
\begin{tabular}{cc}
$\swarrow$ reject unsupported synthesis &
commit accepted parent synthesis $\searrow$
\end{tabular}
\end{tabular}
\caption{Runtime-lifted RLM execution in RLM-FORGE. Ouroboros owns recursive
control and trace state, Hermes executes bounded inner calls, and TraceGuard
decides whether parent synthesis may become committed state.}
\label{fig:architecture}
\end{figure}

\begin{table}[t]
\centering
\caption{RLM advantages preserved by RLM-FORGE at runtime.}
\label{tab:runtime}
\small
\begin{tabular}{%
  >{\raggedright\arraybackslash}p{0.27\linewidth}
  >{\raggedright\arraybackslash}p{0.35\linewidth}
  >{\raggedright\arraybackslash}p{0.27\linewidth}}
\toprule
RLM advantage & Runtime mechanism & Concrete artifact \\
\midrule
Context beyond one prompt & Bounded sub-calls over selected chunks &
\texttt{ooo rlm --truncation-benchmark} \\
Compositional reasoning & Explicit parent-child RLM call graph &
EventStore trace and RLM tree state \\
Intermediate inspectability & Child evidence handles and replayable outputs &
\texttt{benchmarks/*.json} replay \\
Safer synthesis & TraceGuard validates parent claims before mutation &
\texttt{run-traceguard-demo.py} \\
Provider portability & Hermes keeps the inner model swappable &
\texttt{HermesCliRuntime} adapter \\
\bottomrule
\end{tabular}
\end{table}

Table~\ref{tab:runtime} is the main thesis of this paper. The live fixture
does not show a quality delta, but it does show that these RLM control
surfaces exist in a real agent runtime and can be exercised through the
workflow in Figure~\ref{fig:architecture}.

\subsection{Two layers, one orchestration owner}
\label{sec:layers}

The integration has two layers. Ouroboros is the recursion controller and
owns scheduling, termination, guardrails, and trace persistence. Hermes is
the inner LM, called through the existing adapter for one bounded
inference at a time. Hermes never recursively calls Ouroboros or invokes
\texttt{ooo} commands as part of the RLM loop. This constraint is encoded
both in the system prompt sent to Hermes and in the structured contract that
validates its responses.

This ownership split is what lets RLM-FORGE inherit the control advantage of
RLM without pretending Hermes itself has become a new recursive model. Hermes
does local work. Ouroboros owns the recursive state machine. TraceGuard decides
whether a parent is allowed to claim what its children observed.

Each outer-to-inner handoff is a single envelope rendered as the prompt
passed through the \texttt{HermesCliRuntime} adapter. The envelope carries
the following fields.

\begin{itemize}[leftmargin=*]
\item \texttt{mode}. One of \texttt{decompose\_ac}, \texttt{execute\_atomic},
      \texttt{summarize\_chunk}, or \texttt{synthesize\_parent}.
\item \texttt{call\_context}. The values \texttt{call\_id},
      \texttt{parent\_call\_id}, \texttt{rlm\_node\_id}, and
      \texttt{ac\_node\_id}.
\item \texttt{constraints}. The flag \texttt{must\_not\_call\_ouroboros} set
      to \texttt{true}, the ambiguity threshold, and the depth cap.
\item \texttt{context}. Only the chunks, summaries, ancestry, and child
      results selected by the outer scaffold.
\item \texttt{evidence\_handles}. The set of identifiers Hermes is
      permitted to cite.
\end{itemize}

The response carries \texttt{verdict}, \texttt{confidence}, \texttt{result},
evidence references, and residual gaps. All of this is validated before the
outer scaffold mutates state.

\subsection{Sub-call lifecycle}
\label{sec:lifecycle}

For one node at depth $d$ with mode $m$, the lifecycle is as follows.

\begin{enumerate}[leftmargin=*]
\item Validate guardrails ($d \le 5$, ambiguity $\le 0.2$).
\item Bind bounded context and write a call-start trace event with a
      stable \texttt{parent\_call\_id}.
\item Invoke Hermes through \texttt{HermesCliRuntime}.
\item Parse and validate the JSON envelope, including schema version,
      echoed identifiers, and evidence-reference admissibility.
\item Commit accepted results into the AC tree, RLM tree, artifacts, and
      EventStore trace. Invalid responses are recorded as failures and
      do not mutate recursive state.
\item Schedule any follow-up decomposition, atomic execution, summary,
      synthesis, retry, or termination decision.
\end{enumerate}

Hermes may \emph{recommend} that a node is atomic, decomposed, retryable,
or ready for synthesis. Only Ouroboros turns that recommendation into
recursive control flow.

\subsection{Recursive scaffold path}
\label{sec:scaffold}

The first integration path is a dedicated outer-scaffold loop,
\texttt{RLMOuterScaffoldLoop}, in approximately 3{,}200 lines of Python.
It owns RLM node scheduling, AC tree mutations, sub-call dispatch through
the Hermes adapter, parent synthesis, and termination. It is the path
exercised by the user-facing CLI command \texttt{ooo rlm}.

\subsection{Decomposition pipeline path}
\label{sec:decomposition}

The second integration path extends an existing acceptance-criterion
decomposition function, \texttt{decompose\_ac}, with optional parameters
including \texttt{hermes\_runtime}. When that parameter is non-\texttt{None},
the function delegates child-AC generation to a Hermes sub-call. When it is
\texttt{None}, the default path falls through to the original LLM-only flow.
Concretely, the host system commands \texttt{ooo run} and
\texttt{ooo evolve} never pass \texttt{hermes\_runtime}, so they
continue through the tested default path. The same \texttt{HermesCliRuntime}
adapter serves both paths.

This second path is, to our knowledge, an unusual choice. It demonstrates
that an RLM-style boundary can be slotted into an existing
decomposition pipeline incrementally, without replacing the pipeline.
We exercise it through an environment-gated integration test that runs
real Hermes when \texttt{OUROBOROS\_HERMES\_LIVE=1} and is otherwise
skipped to avoid consuming API credits in CI.

\section{Experiment}
\label{sec:experiment}

The experiment is designed to test the runtime claim, not to optimize for a
quality leaderboard. We ask whether the same Hermes-backed stack can execute
a vanilla single call and a recursive RLM call graph, persist both artifacts,
replay the recursive trace, and score omitted-fact safety with a claim-aware
evaluator. The answer is yes. The quality result on the live fixture is a
tie, which is why the runtime lens matters. The demonstrated value is the
control surface that appears around inference, not a score delta on one
synthetic input.

\subsection{Fixture}

We construct a synthetic long-context fixture
\texttt{rlm-long-context-truncation-v1} containing twelve lines split
across six two-line chunks. Four retained facts
\textsc{LC-001}--\textsc{LC-004} fit in the selected chunks and two fact
IDs, \textsc{LC-005} and \textsc{LC-006}, live beyond the retained
truncation boundary.

The fixture is deterministic. Both the vanilla single-call run and the
recursive RLM run receive the same selected-omitted partition and are
scored by the same RLM quality code. The comparison is therefore controlled
across all variables except the inference strategy.

\subsection{Setup}

Hermes is invoked through \texttt{HermesCliRuntime} v0.11.0 against a
custom OpenAI-compatible endpoint. The model and provider are externally
configurable through the Hermes config file. The persisted artifact was
produced with a provider-anonymized frontier-tier chat model routed through
that custom endpoint. The endpoint is private, while public reproduction can
use any Hermes-supported provider. The artifact should be treated as a
systems smoke test rather than a result that generalises across providers.

\begin{table}[t]
\centering
\caption{Environment recorded for the live truncation artifact.}
\label{tab:environment}
\small
\begin{tabular}{ll}
\toprule
Field & Value \\
\midrule
Hermes CLI & v0.11.0 \\
Hermes endpoint & Custom OpenAI-compatible endpoint \\
Inner model & Provider-anonymized frontier-tier chat model \\
Ouroboros dependency & \texttt{5b4bc10e} claim-aware scorer commit \\
RLM-FORGE command & \texttt{ooo rlm --truncation-benchmark} \\
Fixture & \texttt{rlm-long-context-truncation-v1} \\
Sub-calls & vanilla=1, recursive=5 \\
Telemetry not persisted & latency, token count, monetary cost \\
\bottomrule
\end{tabular}
\end{table}

The recursive RLM is configured with \texttt{max\_depth=5} and
\texttt{ambiguity\_threshold=0.2}. The truncation fixture only requires
depth two in practice, with one root sub-call plus four chunk-level sub-calls.

\subsection{Multi-runtime primary portability}

To separate the runtime-contract claim from the older single-provider
artifact, we also build a brownfield live portability harness. The harness
generates an 8-fixture, 3-family, 4-contract matrix plan
(\(8 \times 3 \times 4 = 96\) cells). In this version we execute the
primary contract, \texttt{RLM-FORGE+TraceGuard}, across all 8 fixtures and
all 3 runtime families (\(8 \times 3 = 24\) live cells). Each completed cell
executes four child calls plus one parent synthesis call through the same
\texttt{AgentRuntime} surface, using identical fixtures, prompts, schemas,
scoring, and TraceGuard logic. Only the runtime adapter and model alias vary
by family.

\begin{table}[t]
\centering
\caption{Live primary portability matrix over 8 shared fixtures. This is the
primary RLM-FORGE+TraceGuard matrix, not the full 96-cell baseline sweep.}
\label{tab:portability-primary}
\small
\begin{tabular}{lllll}
\toprule
Family & Runtime & Model alias & Cells & Result \\
\midrule
\texttt{hermes\_glm} & \texttt{HermesCliRuntime} & \texttt{glm-4.7} & 8/8 pass & pass \\
\texttt{claude\_code\_opus47} & \texttt{ClaudeAgentAdapter} & \texttt{opus} & 8/8 pass & pass \\
\texttt{codex\_gpt55} & \texttt{CodexCliRuntime} & \texttt{gpt-5.5} & 8/8 pass & pass \\
\bottomrule
\end{tabular}
\end{table}

The run completes all 24 live cells with 24 mandatory contract passes. The
aggregate status is pass. Hermes+GLM, Claude Code, and Codex each complete
and pass all 8 fixtures under the same RLM-FORGE+TraceGuard contract. The
chunk-only citation trap is still useful as a stress fixture because it
requires parent claims to carry both \texttt{fact\_id} and
\texttt{evidence\_chunk\_id}. In this run every family satisfies that
contract. An earlier Hermes+GLM run exposed the concrete failure mode that
motivates the repair loop. The parent preserved the fact text but omitted the
required evidence handle, causing a \texttt{missing\_evidence\_handle}
rejection. The current implementation now repairs that narrow class by filling
only missing or null evidence-handle fields from the child evidence manifest
and retrying parent synthesis once. If that observed class recurs, the runtime
has a bounded recovery path rather than only a terminal reject-and-log path.
The latest live matrix does not exercise that path because every initial
parent synthesis validates before repair. The repair loop is therefore a
tested runtime-control response, not a claimed benchmark win in the latest
matrix. This does not establish that the full baseline sweep, token usage, or
cost profiles generalise across providers.

\subsection{Results}
\label{sec:results}

\begin{table}[t]
\centering
\caption{Side-by-side comparison on
\texttt{rlm-long-context-truncation-v1}. Both runs use the same
\texttt{HermesCliRuntime}, the same fixture, and the same scoring code.}
\label{tab:results}
\small
\begin{tabular}{lll}
\toprule
Metric & Vanilla single call & Recursive RLM \\
\midrule
Hermes sub-calls & 1 & 5 \\
Quality score & 1.00 & 1.00 \\
Omitted-fact safety & 1.00 & 1.00 \\
Claimed omitted facts & $\emptyset$ & $\emptyset$ \\
Cited retained facts & \textsc{LC-001..004} & \textsc{LC-001..004} \\
Recursive trace score & N/A & 1.00 \\
Parent synthesis observed & N/A & true \\
\bottomrule
\end{tabular}
\end{table}

Table~\ref{tab:results} summarises the corrected result. The recursive RLM
does not outperform the vanilla single call on this fixture. Both runs cite
the same retained facts \textsc{LC-001}--\textsc{LC-004} and avoid positive
claims about omitted facts. The important difference is structural. The
vanilla call returns one final answer, while the recursive run returns a
five-call trace with bounded child contexts, child evidence, and a parent
synthesis step that can be inspected or rejected independently.

\subsection{Omitted-Fact Safety Audit}
\label{sec:mechanism}

The initial artifact reported that the vanilla single-call claimed
\textsc{LC-005} and \textsc{LC-006}. That interpretation was wrong. The
vanilla completion mentioned those IDs inside a residual-gap caveat stating
that omitted facts, if present there, could not be claimed as observed
evidence. A string-matching scorer counted the mention as a positive claim.

The corrected scorer distinguishes structured claim surfaces such as
retained facts, observed facts, evidence references, and unguarded result
claims from non-claim surfaces such as residual gaps and truncation reports.
Under that scoring rule, both vanilla and recursive RLM return an empty
claimed-omitted-facts list.

The recursive design still provides the key runtime boundary. Child sub-calls
are supplied bounded chunks, parent synthesis consumes accepted child results,
and evidence handles can be validated before state mutation.
This changes the error geometry. Unsupported claims are no longer just text
inside a final answer. They become contract violations at a named synthesis
boundary. A model can still fabricate unsupported IDs, but the runtime now
has a precise place to reject them.

\subsection{TraceGuard Parent Synthesis Enforcement}

TraceGuard is the enforcement layer for that contract. It consumes an
accepted child evidence manifest and a parent synthesis object. Claims in
\texttt{retained\_facts}, \texttt{observed\_facts}, and
\texttt{evidence\_references} are accepted only when their fact ID and
evidence handle match the manifest. Omitted facts are rejected with an
\texttt{unsupported\_fact\_id} reason. Chunk handles without fact IDs are
rejected with \texttt{chunk\_handle\_without\_fact}. This makes the
evidence boundary executable rather than merely descriptive.

\paragraph{Acceptance property.}
Let $M \subseteq F \times H$ be the accepted child evidence manifest, where
$F$ is the set of fact identifiers and $H$ is the set of evidence handles.
Let $\mathrm{Claims}(S)$ be the structured fact claims extracted from a
parent synthesis $S$. TraceGuard accepts $S$ only if
\[
  \forall (f,h) \in \mathrm{Claims}(S), \quad (f,h) \in M.
\]
The implementation builds $M$ only from accepted child outputs, enumerates
every structured claim surface in $S$, and rejects on the first unknown fact
identifier, omitted fact identifier, chunk handle without a fact, or
fact-handle pair absent from $M$. Acceptance is returned only after this scan
completes. This is a provenance property relative to the manifest, not a
semantic proof that the underlying natural-language fact is true.

The public interface is shown below.

\begin{lstlisting}[language=Python]
validate_parent_synthesis(
    evidence_manifest=build_manifest_from_fixture(fixture),
    parent_synthesis=parent_json,
)
\end{lstlisting}

The no-API demo accepts safe parent synthesis and rejects an omitted fact.
This separates score parity on the live fixture from the new capability,
which is enforceable provenance support at the recursive synthesis boundary.

\subsection{Reproduction}

The benchmark can be reproduced live with one command.

\begin{lstlisting}[language=bash]
ooo rlm --truncation-benchmark
\end{lstlisting}

Without an API key, the persisted artifact can be replayed.

\begin{lstlisting}[language=bash]
python -m rlm_forge.replay \
    benchmarks/rlm-long-context-truncation-v1.json
\end{lstlisting}

Both commands print the headline metrics from
Table~\ref{tab:results}.

The TraceGuard enforcement demo runs without Hermes.

\begin{lstlisting}[language=bash]
python3 scripts/run-traceguard-demo.py
\end{lstlisting}

It accepts a safe parent synthesis, rejects an omitted-fact claim, and
rejects chunk-only evidence that does not identify supported fact IDs.

The 24-cell live primary portability matrix can be rerun with provider
credentials configured.

\begin{lstlisting}[language=bash]
HERMES_INFERENCE_PROVIDER=zai \
GLM_API_KEY=... \
uv run --extra dev --extra live \
  python scripts/run-live-portability-matrix.py \
  --mode live-primary --fixtures 8 --families all \
  --timeout-seconds 240 \
  --output-prefix live-portability-primary
\end{lstlisting}

We additionally include a Hermes-free scorer micro-suite.

\begin{lstlisting}[language=bash]
python3 scripts/run-claim-aware-suite.py
\end{lstlisting}

The current suite contains seven controlled completion shapes and all seven
pass under the corrected scorer. The cases cover guarded omitted-ID mentions,
positive omitted-fact assertions, omitted evidence references, chunk-only
citations without fact evidence, missing truncation-boundary reports, and
missing retained facts.

We also include a generated synthetic scorer benchmark.

\begin{lstlisting}[language=bash]
python3 scripts/run-synthetic-omitted-fact-benchmark.py
\end{lstlisting}

It creates 108 truncation fixtures by varying fact count, retained/omitted
ratio, distractor density, and omitted-claim target. Each fixture is scored
against seven deterministic completion strategies, producing 756 scorer
sanity checks. All pass in the committed artifact. This benchmark validates
the measurement harness, not live-model quality.

Finally, we include an unsupported-claim-rate contract ablation.

\begin{lstlisting}[language=bash]
python3 scripts/run-unsupported-claim-rate-benchmark.py
\end{lstlisting}

This benchmark creates 72 generated fixtures and evaluates six execution
contracts. They are loose single call, guarded single call, chunk-only
map-reduce, leaky map-reduce, evidence-gated Hermes-RLM, and
Hermes-RLM-shaped recursion without evidence gating. The evidence-gated
Hermes-RLM contract has an unsupported-claim rate of 0.0. The ungated
RLM-shaped contract has an unsupported-claim rate of 1.0. The result is
deterministic and should be read as a contract ablation, not as a live-model
comparison.

\section{Discussion}
\label{sec:discussion}

RLM-FORGE is best understood as a systems argument. A model-only reading of
RLM asks whether a recursive call beats a vanilla call on a benchmark. That
question is important, but it is too narrow for systems work. A runtime
reading asks whether the RLM pattern exposes control surfaces that ordinary
agent infrastructure can use immediately.

\paragraph{Contribution under score parity.}
The corrected live fixture reports score parity. This does not negate the
runtime claim. A vanilla call can answer a controlled truncation fixture
correctly. RLM-FORGE is not claiming otherwise. The contribution is that the
recursive run produces a different artifact, namely a traceable computation with
bounded child calls, explicit evidence handles, and a parent synthesis gate.
That artifact is easier to debug, replay, audit, and constrain than a single
final completion. In deployed agent systems, those properties are often
valuable independently of benchmark score.

\paragraph{Runtime RLM changes the failure mode.}
Without evidence gating, an RLM-shaped policy can hallucinate just as badly
as a loose single call. The contract ablation is intentionally blunt. Ungated
recursion fails at a 1.0 unsupported-claim rate, while the evidence-gated
Hermes-RLM contract is 0.0 on the deterministic suite. The
lesson is not that recursion magically prevents unsupported claims. The
lesson is that recursion creates a location where unsupported claims can be
detected and rejected before they become committed state.

\paragraph{The platform effect is the real upside.}
Because Hermes keeps the inner model provider-swappable, RLM-FORGE turns
recursive inference into a platform feature rather than a checkpoint feature.
The same outer scaffold can call different inner models, preserve the same
trace schema, and enforce the same parent-synthesis contract. This is the
practical route by which RLM can spread. It does not require every provider to
ship a native recursive model. It makes recursive execution a capability of
the runtime that sits above providers.

\paragraph{Memory-shaped recursion.}
RLM-FORGE currently treats each run as an isolated recursive execution.
Hermes makes a stronger future direction possible. Decomposition traces,
provider-specific failures, schema repairs, and TraceGuard rejections can
become persistent operational priors for later recursive runs. Memory is not
evidence. Memory is a prior over how to ask for evidence. Fresh parent claims
must still cite fresh child evidence and pass TraceGuard before state mutation.

\paragraph{Integration scope.}
We added no new sandbox, REPL, or transport. The integration that matters
lives in approximately 4{,}500 lines of Python (recursive scaffold and
decomposition extension, excluding tests and concept documentation), all
calling through one pre-existing \texttt{HermesCliRuntime} adapter.
Researchers and practitioners with similar agent runtimes should be able to
reproduce the integration in comparable time.

\section{Future Work}
\label{sec:future-memory}

Hermes does not train the underlying model during RLM execution. The future
opportunity is different. Hermes memory, skills, and traces can give the
recursive scaffold a persistent operational memory. Prior decompositions,
provider-specific failure modes, schema repairs, and TraceGuard rejection
patterns can be recalled as runtime priors for later recursive runs.

The key separation is this.

\begin{quote}
Memory is not evidence. Memory is a prior over how to ask for evidence.
\end{quote}

In this design, memory should not be allowed to support a factual parent
claim. Only fresh child evidence from the current run can do that. Memory can
shape the policy for the next run. It can influence how to decompose a task
family, which provider to route extraction or synthesis to, which schema to
require, and which retry prompt to use after a guardrail rejection. A future
memory-shaped run could recall missing-handle rejections as schema-repair
priors. It could require child outputs to echo both \texttt{fact\_id} and
\texttt{evidence\_chunk\_id}, and retry a child call when a handle is null.
The repaired parent would still need to cite fresh evidence from the new trace
before TraceGuard accepts it.

This direction is useful only if memory is curated. Raw progress logs,
provider chatter, and fixture answers should not become durable recall. The
memory record should be an abstract operational signal such as "this provider
is slow but usually schema-stable" or "this task family needs fact-level
handles rather than chunk-only handles." Benchmark use also requires
contamination controls. These include no-memory baselines, answer-memory
disabled, fixture content excluded from durable memory, and separate
trace-memory policy runs.

\section{Limitations}
\label{sec:limits}

We are explicit about what this paper does and does not show.

\begin{enumerate}[leftmargin=*]
\item \textbf{One score-bearing live truncation fixture.} The corrected
Hermes truncation artifact is a tie on one truncation regime with a
particular ratio of selected to omitted chunks. The live primary matrix
expands contract coverage across 8 fixtures, but it is not a broad
model-quality benchmark.
\item \textbf{Primary-only provider coverage.} We now run the primary
RLM-FORGE+TraceGuard contract across 8 fixtures, Hermes+GLM, Claude Code,
and Codex CLI. This is stronger than a smoke test, but it is still not the
full 96-cell baseline sweep. The primary aggregate is pass. The unanswered
question is whether the secondary variants show the same pattern when executed
across all providers.
\item \textbf{Incomplete cost and token telemetry.} The portability run
records wall-clock latency, but not token counts or monetary cost. Those
fields should be captured in a future full multi-provider run.
\item \textbf{No comparison to RLM-Qwen3-8B.} The natively recursive
model post-trained by \citet{zhang2025rlm} is not in our comparison set.
\item \textbf{No new RLM theory.} The recursive paradigm is the
contribution of \citet{zhang2025rlm}. We contribute a runtime lifting
recipe, an evidence-gated synthesis contract, and a specific scoring axis.
\item \textbf{Synthetic fact identifiers.} Our omitted-fact-safety
measurement depends on facts that the fixture explicitly labels with
identifiers. Generalising the metric to natural-language settings
requires a robust fact extractor.
\end{enumerate}

\section{Related Work}
\label{sec:related}

\citet{zhang2025rlm} introduce RLM as an inference strategy and evaluate
it across deep research, information aggregation, code repository
understanding, and a synthetic pairwise reasoning task. We treat their
paradigm as foundational and contribute an implementation note plus a
focused evaluation axis.

The general space of long-context handling for LLMs is broad. Retrieval
agents, summary agents, and code-execution agents such as CodeAct
~\citep{wang2024codeact} all attempt to manage long inputs without
explicit recursion. RLM differs from these by treating recursive
sub-calls as first-class operations on the input rather than as
auxiliary tools.

The Hermes Agent runtime~\citep{nous2025hermes} is part of a recent
generation of agent runtimes that treat the LLM as a callable function
with structured I/O, alongside related systems aimed at production
agent deployment. Our integration takes advantage of this design without
modifying it.

\section{Conclusion}
\label{sec:conclusion}

We have presented RLM-FORGE, an open-source runtime lifting of Recursive
Language Models onto the Hermes Agent runtime. The central systems claim is
that the operational advantages of RLM do not require a new model architecture
to become useful. They can be exposed as runtime properties, including bounded
sub-calls, recursive trace state, replayable child evidence,
provider-swappable inner execution, and evidence-gated parent synthesis. The
committed live fixture shows score parity after claim-aware
rescoring, but the recursive run produces a computation the runtime can
inspect and enforce. The 24-cell live primary matrix further shows that the
same evidence-gated runtime contract can pass across Hermes+GLM, Claude Code,
and Codex. That is the systems contribution. The natural next step is
memory-shaped recursion, where prior traces improve decomposition, routing,
and repair policy while fresh child evidence remains the only basis for parent
claims. We offer the code,
fixtures, and persisted artifacts under an MIT licence at
\url{https://github.com/Q00/rlm-forge}.

\section*{Reproducibility}

All code, fixtures, persisted side-by-side artifacts, replay tooling,
and configuration files are available at
\url{https://github.com/Q00/rlm-forge}. The Ouroboros branch
hosting the upstream RLM modules is at
\url{https://github.com/Q00/ouroboros/tree/ooo/run/seed_2d9f9ccfa5aa}.
The implementation artifact depends on the Ouroboros branch by git
reference until the upstream release lands on PyPI.

\bibliographystyle{plainnat}
\bibliography{refs}

\end{document}
